% !TEX root = chapter-standalone.tex



\section{New functions from old}
\label{sec:new-functions-from-old}

Analogous to our discussion of how to build new sets from old, we discuss here some basic ways to build new functions from old. 

\subsection{Products}

\begin{ctdefinition}[Product of functions]
    \label{def:product-of-functions}
    \SYNDEF{product of functions}
    Given functions~$\mapa\colon \setA \sto \setB$ and~$\mapb\colon \setC \sto \setD$, their \emph{product} is the function
    \begin{equation}
        \defmapperiod{
            \mapa \funcprod \mapb
        }{
            \setA \cartprod \setC
        }{
            \sto
        }{
            \setB\cartprod \setD
        }{
            \tup{\setAel, \setCel}
        }{
            \tup{\mapa(\setAel),\mapb(\setCel)}
        }
    \end{equation}
\end{ctdefinition}

\begin{example}
    Consider the functions
    \begin{equation}
        \defmapcomma{
            \mapa
        }{
            \reals
        }{
            \sto
        }{
            \reals
        }{
            \setAel
        }{
            \setAel^2
        }
    \quad \quad \text{ and }\quad \quad
        \defmapperiod{
            \mapb
        }{
            \reals
        }{
            \sto
        }{
            \reals
        }{
            \setBel
        }{
            \setBel+1
        }
    \end{equation}
Their product is 
\begin{equation}
        \defmapperiod{
            \mapa \funcprod \mapb 
        }{
            \reals \cartprod \reals
        }{
            \sto
        }{
            \reals \cartprod \reals
        }{
            \tup{\setAel,\setBel}
        }{
            \tup{\setAel^2,\setBel+1}
        }
\end{equation}
    So, for example, 
\begin{equation}
    (\mapa \funcprod \mapb)(\tup{2,5})=\tup{2^2,5+1}=\tup{4,6}.
\end{equation}
\end{example}



\subsection{Sums}

\begin{ctdefinition}[Sum of functions]
    \label{def:sum-of-functions}
    \SYNDEF{sum of functions}
    Given functions~$\mapa\colon \setA \sto \setB$ and~$\mapb\colon \setC \sto \setD$, their \emph{sum} is the function
    \begin{equation}
        \begin{aligned}
            \mapa \funcsum \mapb \colon \setA \setdisunion \setC & \sto \setB\setdisunion \setD \\
            \disunionA{\setAel}                                  & \mapsto \disunionA{\mapa(\setAel)}, \\
            \disunionB{\setCel}                                  & \mapsto \disunionB{\mapb(\setCel)}.
            \\
        \end{aligned}
    \end{equation}
\end{ctdefinition}

\begin{example}
    Consider the functions
    \begin{equation}
        \defmapcomma{
            \mapa
        }{
            \wnumbers
        }{
            \sto
        }{
            \natnumbers
        }{
            \setAel
        }{
            \setAel^2
        }
\quad \quad  \text{ and } \quad \quad 
        \defmapperiod{
            \mapb
        }{
            \reals
        }{
            \sto
        }{
            \reals
        }{
            \setBel
        }{
            \setBel^3
        }
    \end{equation}
    Their sum is 
    \begin{equation}
        \begin{aligned}
            \mapa \funcsum \mapb \colon \wnumbers \setdisunion \reals & \sto \natnumbers \setdisunion \reals \\
            \disunionA{\setAel}                                       & \mapsto \disunionA{\setAel^2}, \\
            \disunionB{\setBel}                                       & \mapsto \disunionB{\setBel^3}.
            \\
        \end{aligned}
    \end{equation}
    So, for example,
    \begin{equation}
    (\mapa \funcsum \mapb)(\tup{1,3})=\tup{1,\mapa(3)}=\tup{1,3^2}=\tup{1,9}
    \end{equation}
    and
    \begin{equation}
        (\mapa \funcsum \mapb)(\tup{2,-3})=\tup{2,\mapb(-3)}=\tup{2,(-3)^3}=\tup{2,-27}.
    \end{equation}
\end{example}



\devel{
\subsection{exponentiation}



\begin{ctdefinition}[Exponentiation of functions]
    \label{def:exponentiation-of-functions}
    \SYNDEF{exponentiation of functions}
    Given a function~$\mapa\colon \setA \sto \setB$ and a set $\setC$, the \emph{exponentiation} of $\mapa$, with base $\setC$, is the function
    \begin{equation}
        \begin{aligned}
            \setC^\mapa \colon \setC^\setB & \sto \setC^\setA \\
            \varphi                        & \mapsto \mapa \mthen \varphi.
        \end{aligned}
    \end{equation}
\end{ctdefinition}

\begin{remark}
    Note that under exponentiation, the order in which $\setA$ and $\setB$ appear becomes ``switched''.
\end{remark}

\todotext{Make at least two examples; a more general one, and one involving the power set (base 2).
    Oh, actually the latter must maybe come later, after we've explained the alternative function way of describing the powerset.
    ..}



}